% Erratum E14 for inclusion in main.tex.
% Suggested use: % Erratum E14 for inclusion in main.tex.
% Suggested use: % Erratum E14 for inclusion in main.tex.
% Suggested use: % Erratum E14 for inclusion in main.tex.
% Suggested use: \input{erratum_E14_landweber_argmax_domain}

\subsection*{Erratum E14: the maximization domain in the proof of Proposition~3.4}
\label{s:erratum:E14}

In the proof of Proposition~3.4, the line preceding the
normalization $x_0=0$ and $\sigma=1$ should not maximize over all
$\lambda\geq0$.  The correct statement is that
\begin{align}
\label{e:erratum:E14:corrected-argmax}
        \varepsilon_k
        =
        \argmax_{0\leq \lambda\leq 1/\sigma}
        \lambda^{2(r+\nu)}(1-\sigma\lambda)^{2k},
        \qquad k\geq0 .
\end{align}
Equivalently, after the normalization $\sigma=1$, the maximization is over
$0\leq\lambda\leq1$.  This is the relevant interval, since the step-size
condition
\[
        0 < \sigma \norm{T^* T}{} \leq 1
\]
implies
\[
        \operatorname{supp}\mu_{x^\dag-x_0}
        \subset [0,\|T^*T\|]
        \subset [0,1/\sigma].
\]
The formulation with $\lambda\geq0$ is false: outside the stability
interval $[0,1/\sigma]$ the factor $(1-\sigma\lambda)^{2k}$ grows again,
and the displayed function is not maximized at $\varepsilon_k$.

We verify \eqref{e:erratum:E14:corrected-argmax}.  Put
\[
        a:=r+\nu>0,
        \qquad
        f_k(\lambda)
        :=
        \lambda^{2a}(1-\sigma\lambda)^{2k} .
\]
For $k=0$, $f_0(\lambda)=\lambda^{2a}$ is maximized on
$[0,1/\sigma]$ at $\lambda=1/\sigma$, which is precisely
\[
        \varepsilon_0
        =
        \frac{1}{\sigma}\frac{a}{a}
        =
        \frac{1}{\sigma} .
\]
For $k\geq1$, consider $0<\lambda<1/\sigma$.  Then
\[
        \frac{d}{d\lambda}\log f_k(\lambda)
        =
        \frac{2a}{\lambda}
        -
        \frac{2k\sigma}{1-\sigma\lambda} .
\]
The unique critical point is therefore determined by
\[
        a(1-\sigma\lambda)=k\sigma\lambda,
\]
that is,
\[
        \lambda
        =
        \frac{1}{\sigma}\frac{a}{k+a}
        =
        \frac{1}{\sigma}\frac{r+\nu}{k+r+\nu}
        =
        \varepsilon_k .
\]
Moreover,
\[
        \frac{d^2}{d\lambda^2}\log f_k(\lambda)
        =
        -\frac{2a}{\lambda^2}
        -
        \frac{2k\sigma^2}{(1-\sigma\lambda)^2}
        <0,
\]
so this critical point is the unique maximum in the open interval.  Since
$f_k(0)=f_k(1/\sigma)=0$ for $k\geq1$, it is also the maximum on
$[0,1/\sigma]$.

This correction does not affect the proof of Proposition~3.4.  All
spectral integrals used there are supported in
$[0,\|T^*T\|]\subset[0,1/\sigma]$, and the subsequent estimates use
only this stability interval.


\subsection*{Erratum E14: the maximization domain in the proof of Proposition~3.4}
\label{s:erratum:E14}

In the proof of Proposition~3.4, the line preceding the
normalization $x_0=0$ and $\sigma=1$ should not maximize over all
$\lambda\geq0$.  The correct statement is that
\begin{align}
\label{e:erratum:E14:corrected-argmax}
        \varepsilon_k
        =
        \argmax_{0\leq \lambda\leq 1/\sigma}
        \lambda^{2(r+\nu)}(1-\sigma\lambda)^{2k},
        \qquad k\geq0 .
\end{align}
Equivalently, after the normalization $\sigma=1$, the maximization is over
$0\leq\lambda\leq1$.  This is the relevant interval, since the step-size
condition
\[
        0 < \sigma \norm{T^* T}{} \leq 1
\]
implies
\[
        \operatorname{supp}\mu_{x^\dag-x_0}
        \subset [0,\|T^*T\|]
        \subset [0,1/\sigma].
\]
The formulation with $\lambda\geq0$ is false: outside the stability
interval $[0,1/\sigma]$ the factor $(1-\sigma\lambda)^{2k}$ grows again,
and the displayed function is not maximized at $\varepsilon_k$.

We verify \eqref{e:erratum:E14:corrected-argmax}.  Put
\[
        a:=r+\nu>0,
        \qquad
        f_k(\lambda)
        :=
        \lambda^{2a}(1-\sigma\lambda)^{2k} .
\]
For $k=0$, $f_0(\lambda)=\lambda^{2a}$ is maximized on
$[0,1/\sigma]$ at $\lambda=1/\sigma$, which is precisely
\[
        \varepsilon_0
        =
        \frac{1}{\sigma}\frac{a}{a}
        =
        \frac{1}{\sigma} .
\]
For $k\geq1$, consider $0<\lambda<1/\sigma$.  Then
\[
        \frac{d}{d\lambda}\log f_k(\lambda)
        =
        \frac{2a}{\lambda}
        -
        \frac{2k\sigma}{1-\sigma\lambda} .
\]
The unique critical point is therefore determined by
\[
        a(1-\sigma\lambda)=k\sigma\lambda,
\]
that is,
\[
        \lambda
        =
        \frac{1}{\sigma}\frac{a}{k+a}
        =
        \frac{1}{\sigma}\frac{r+\nu}{k+r+\nu}
        =
        \varepsilon_k .
\]
Moreover,
\[
        \frac{d^2}{d\lambda^2}\log f_k(\lambda)
        =
        -\frac{2a}{\lambda^2}
        -
        \frac{2k\sigma^2}{(1-\sigma\lambda)^2}
        <0,
\]
so this critical point is the unique maximum in the open interval.  Since
$f_k(0)=f_k(1/\sigma)=0$ for $k\geq1$, it is also the maximum on
$[0,1/\sigma]$.

This correction does not affect the proof of Proposition~3.4.  All
spectral integrals used there are supported in
$[0,\|T^*T\|]\subset[0,1/\sigma]$, and the subsequent estimates use
only this stability interval.


\subsection*{Erratum E14: the maximization domain in the proof of Proposition~3.4}
\label{s:erratum:E14}

In the proof of Proposition~3.4, the line preceding the
normalization $x_0=0$ and $\sigma=1$ should not maximize over all
$\lambda\geq0$.  The correct statement is that
\begin{align}
\label{e:erratum:E14:corrected-argmax}
        \varepsilon_k
        =
        \argmax_{0\leq \lambda\leq 1/\sigma}
        \lambda^{2(r+\nu)}(1-\sigma\lambda)^{2k},
        \qquad k\geq0 .
\end{align}
Equivalently, after the normalization $\sigma=1$, the maximization is over
$0\leq\lambda\leq1$.  This is the relevant interval, since the step-size
condition
\[
        0 < \sigma \norm{T^* T}{} \leq 1
\]
implies
\[
        \operatorname{supp}\mu_{x^\dag-x_0}
        \subset [0,\|T^*T\|]
        \subset [0,1/\sigma].
\]
The formulation with $\lambda\geq0$ is false: outside the stability
interval $[0,1/\sigma]$ the factor $(1-\sigma\lambda)^{2k}$ grows again,
and the displayed function is not maximized at $\varepsilon_k$.

We verify \eqref{e:erratum:E14:corrected-argmax}.  Put
\[
        a:=r+\nu>0,
        \qquad
        f_k(\lambda)
        :=
        \lambda^{2a}(1-\sigma\lambda)^{2k} .
\]
For $k=0$, $f_0(\lambda)=\lambda^{2a}$ is maximized on
$[0,1/\sigma]$ at $\lambda=1/\sigma$, which is precisely
\[
        \varepsilon_0
        =
        \frac{1}{\sigma}\frac{a}{a}
        =
        \frac{1}{\sigma} .
\]
For $k\geq1$, consider $0<\lambda<1/\sigma$.  Then
\[
        \frac{d}{d\lambda}\log f_k(\lambda)
        =
        \frac{2a}{\lambda}
        -
        \frac{2k\sigma}{1-\sigma\lambda} .
\]
The unique critical point is therefore determined by
\[
        a(1-\sigma\lambda)=k\sigma\lambda,
\]
that is,
\[
        \lambda
        =
        \frac{1}{\sigma}\frac{a}{k+a}
        =
        \frac{1}{\sigma}\frac{r+\nu}{k+r+\nu}
        =
        \varepsilon_k .
\]
Moreover,
\[
        \frac{d^2}{d\lambda^2}\log f_k(\lambda)
        =
        -\frac{2a}{\lambda^2}
        -
        \frac{2k\sigma^2}{(1-\sigma\lambda)^2}
        <0,
\]
so this critical point is the unique maximum in the open interval.  Since
$f_k(0)=f_k(1/\sigma)=0$ for $k\geq1$, it is also the maximum on
$[0,1/\sigma]$.

This correction does not affect the proof of Proposition~3.4.  All
spectral integrals used there are supported in
$[0,\|T^*T\|]\subset[0,1/\sigma]$, and the subsequent estimates use
only this stability interval.


\subsection*{Erratum E14: the maximization domain in the proof of Proposition~3.4}
\label{s:erratum:E14}

In the proof of Proposition~3.4, the line preceding the
normalization $x_0=0$ and $\sigma=1$ should not maximize over all
$\lambda\geq0$.  The correct statement is that
\begin{align}
\label{e:erratum:E14:corrected-argmax}
        \varepsilon_k
        =
        \argmax_{0\leq \lambda\leq 1/\sigma}
        \lambda^{2(r+\nu)}(1-\sigma\lambda)^{2k},
        \qquad k\geq0 .
\end{align}
Equivalently, after the normalization $\sigma=1$, the maximization is over
$0\leq\lambda\leq1$.  This is the relevant interval, since the step-size
condition
\[
        0 < \sigma \norm{T^* T}{} \leq 1
\]
implies
\[
        \operatorname{supp}\mu_{x^\dag-x_0}
        \subset [0,\|T^*T\|]
        \subset [0,1/\sigma].
\]
The formulation with $\lambda\geq0$ is false: outside the stability
interval $[0,1/\sigma]$ the factor $(1-\sigma\lambda)^{2k}$ grows again,
and the displayed function is not maximized at $\varepsilon_k$.

We verify \eqref{e:erratum:E14:corrected-argmax}.  Put
\[
        a:=r+\nu>0,
        \qquad
        f_k(\lambda)
        :=
        \lambda^{2a}(1-\sigma\lambda)^{2k} .
\]
For $k=0$, $f_0(\lambda)=\lambda^{2a}$ is maximized on
$[0,1/\sigma]$ at $\lambda=1/\sigma$, which is precisely
\[
        \varepsilon_0
        =
        \frac{1}{\sigma}\frac{a}{a}
        =
        \frac{1}{\sigma} .
\]
For $k\geq1$, consider $0<\lambda<1/\sigma$.  Then
\[
        \frac{d}{d\lambda}\log f_k(\lambda)
        =
        \frac{2a}{\lambda}
        -
        \frac{2k\sigma}{1-\sigma\lambda} .
\]
The unique critical point is therefore determined by
\[
        a(1-\sigma\lambda)=k\sigma\lambda,
\]
that is,
\[
        \lambda
        =
        \frac{1}{\sigma}\frac{a}{k+a}
        =
        \frac{1}{\sigma}\frac{r+\nu}{k+r+\nu}
        =
        \varepsilon_k .
\]
Moreover,
\[
        \frac{d^2}{d\lambda^2}\log f_k(\lambda)
        =
        -\frac{2a}{\lambda^2}
        -
        \frac{2k\sigma^2}{(1-\sigma\lambda)^2}
        <0,
\]
so this critical point is the unique maximum in the open interval.  Since
$f_k(0)=f_k(1/\sigma)=0$ for $k\geq1$, it is also the maximum on
$[0,1/\sigma]$.

This correction does not affect the proof of Proposition~3.4.  All
spectral integrals used there are supported in
$[0,\|T^*T\|]\subset[0,1/\sigma]$, and the subsequent estimates use
only this stability interval.
